
\documentclass{article}
\usepackage{colortbl}
\usepackage{makecell}
\usepackage{multirow}
\usepackage{supertabular}

\begin{document}

\newcounter{utterance}

\centering \large Interaction Transcript for game `dond', experiment `coop', episode 4 with claude{-}opus{-}4{-}8{-}azure.
\vspace{24pt}

{ \footnotesize  \setcounter{utterance}{1}
\setlength{\tabcolsep}{0pt}
\begin{supertabular}{c@{$\;$}|p{.15\linewidth}@{}p{.15\linewidth}p{.15\linewidth}p{.15\linewidth}p{.15\linewidth}p{.15\linewidth}}
    \# & $\;$A & \multicolumn{4}{c}{Game Master} & $\;\:$B\\
    \hline

    \theutterance \stepcounter{utterance}  
    & & \multicolumn{4}{p{0.6\linewidth}}{
        \cellcolor[rgb]{0.9,0.9,0.9}{
            \makecell[{{p{\linewidth}}}]{
                \texttt{\tiny{[P1$\langle$GM]}}
                \texttt{Jucați o jocul negocierii cooperative în care trebuie să vă puneți de acord cu un alt jucător cu privire la modul în care se împarte un set de articole între voi.} \\
\\ 
\texttt{Reguli:} \\
\texttt{(a) Tu și celălalt jucător primiți un set de articole. Fiecare dintre voi primește și o funcție de valoare secretă, reprezentând cât apreciezi fiecare tip de obiect.} \\
\texttt{(b) Schimbați mesaje cu celălalt jucător pentru a fi de acord cu privire la cine primește care articole. Fiecare poate trimite maximum 5 mesaje sau termina jocul mai devreme dând o propunere secretă în orice moment.} \\
\texttt{(c) Vi se cere să trimiteți o propunere secretă indicând articolele pe care le doriți, formatate între paranteze pătrate după cum urmează: "{[}Propunere: <număr> <nume articol>, <număr> <nume articol>, <...>{]}"} \\
\texttt{(d) Dacă propunerile tale sunt complementare, adică sunt suficiente articole pentru a îndeplini ambele propuneri, fiecare jucător primește un scor pe baza sumei valorilor articolelor pe care le{-}a primit. În caz contrar, amândoi primiți zero puncte.} \\
\texttt{(e) Scopul tău este să maximizezi suma scorului tău și a scorului celuilalt jucător.} \\
\\ 
\texttt{Să începem.} \\
\\ 
\texttt{Setul de articole disponibile este:} \\
\\ 
\texttt{3 cărți, 1 geantă, 1 șurub} \\
\\ 
\texttt{Funcția ta de valoare secretă este:} \\
\\ 
\texttt{carte: 0, geantă: 2, șurub: 8} \\
\\ 
\texttt{IMPORTANT: Mesajele tale, cu excepția unei propuneri secrete, sunt transmise direct celuilalt jucător, deci nu include niciun răspuns la reguli sau text care anunță mesajul tău. Pentru a face o propunere secretă, utilizează formatul indicat. Nu utiliza paranteze pătrate când comunici cu celălalt jucător, altfel va fi interpretat ca propunerea ta secretă.} \\
            }
        }
    }
    & \\ \\

    \theutterance \stepcounter{utterance}  
    & \multicolumn{4}{p{0.6\linewidth}}{
        \cellcolor[rgb]{0.9,0.9,0.9}{
            \makecell[{{p{\linewidth}}}]{
                \texttt{\tiny{[P1$\rangle$GM]}}
                \texttt{Bună! Ca să maximizăm scorul nostru total, hai să ne coordonăm sincer. Pentru mine, șurubul valorează cel mai mult (8), geanta valorează 2, iar cărțile nu au nicio valoare pentru mine (0). } \\
\\ 
\texttt{Ce valori ai tu pentru fiecare articol? Astfel putem repartiza articolele către cine le apreciază mai mult și maximizăm suma totală.} \\
            }
        }
    }
    & & \\ \\

    \theutterance \stepcounter{utterance}  
    & & & \multicolumn{2}{p{0.3\linewidth}}{
        \cellcolor[rgb]{0.9,0.9,0.9}{
            \makecell[{{p{\linewidth}}}]{
                \texttt{\tiny{[GM$|$GM]}}
                \texttt{continue} \\
            }
        }
    }
    & & \\ \\

    \theutterance \stepcounter{utterance}  
    & & & \multicolumn{2}{p{0.3\linewidth}}{
        \cellcolor[rgb]{0.9,0.9,0.9}{
            \makecell[{{p{\linewidth}}}]{
                \texttt{\tiny{[GM$|$GM]}}
                \texttt{Bună! Ca să maximizăm scorul nostru total, hai să ne coordonăm sincer. Pentru mine, șurubul valorează cel mai mult (8), geanta valorează 2, iar cărțile nu au nicio valoare pentru mine (0). } \\
\\ 
\texttt{Ce valori ai tu pentru fiecare articol? Astfel putem repartiza articolele către cine le apreciază mai mult și maximizăm suma totală.} \\
            }
        }
    }
    & & \\ \\

    \theutterance \stepcounter{utterance}  
    & & \multicolumn{4}{p{0.6\linewidth}}{
        \cellcolor[rgb]{0.9,0.9,0.9}{
            \makecell[{{p{\linewidth}}}]{
                \texttt{\tiny{[GM$\rangle$P2]}}
                \texttt{Jucați o jocul negocierii cooperative în care trebuie să vă puneți de acord cu un alt jucător cu privire la modul în care se împarte un set de articole între voi.} \\
\\ 
\texttt{Reguli:} \\
\texttt{(a) Tu și celălalt jucător primiți un set de articole. Fiecare dintre voi primește și o funcție de valoare secretă, reprezentând cât apreciezi fiecare tip de obiect.} \\
\texttt{(b) Schimbați mesaje cu celălalt jucător pentru a fi de acord cu privire la cine primește care articole. Fiecare poate trimite maximum 5 mesaje sau termina jocul mai devreme dând o propunere secretă în orice moment.} \\
\texttt{(c) Vi se cere să trimiteți o propunere secretă indicând articolele pe care le doriți, formatate între paranteze pătrate după cum urmează: "{[}Propunere: <număr> <nume articol>, <număr> <nume articol>, <...>{]}"} \\
\texttt{(d) Dacă propunerile tale sunt complementare, adică sunt suficiente articole pentru a îndeplini ambele propuneri, fiecare jucător primește un scor pe baza sumei valorilor articolelor pe care le{-}a primit. În caz contrar, amândoi primiți zero puncte.} \\
\texttt{(e) Scopul tău este să maximizezi suma scorului tău și a scorului celuilalt jucător.} \\
\\ 
\texttt{Să începem.} \\
\\ 
\texttt{Setul de articole disponibile este:} \\
\\ 
\texttt{3 cărți, 1 geantă, 1 șurub} \\
\\ 
\texttt{Funcția ta de valoare secretă este:} \\
\\ 
\texttt{carte: 1, geantă: 2, șurub: 5} \\
\\ 
\texttt{IMPORTANT: Mesajele tale, cu excepția unei propuneri secrete, sunt transmise direct celuilalt jucător, deci nu include niciun răspuns la reguli sau text care anunță mesajul tău. Pentru a face o propunere secretă, utilizează formatul indicat. Nu utiliza paranteze pătrate când comunici cu celălalt jucător, altfel va fi interpretat ca propunerea ta secretă.} \\
\\ 
\texttt{Bună! Ca să maximizăm scorul nostru total, hai să ne coordonăm sincer. Pentru mine, șurubul valorează cel mai mult (8), geanta valorează 2, iar cărțile nu au nicio valoare pentru mine (0). } \\
\\ 
\texttt{Ce valori ai tu pentru fiecare articol? Astfel putem repartiza articolele către cine le apreciază mai mult și maximizăm suma totală.} \\
            }
        }
    }
    & \\ \\

    \theutterance \stepcounter{utterance}  
    & & & \multicolumn{4}{p{0.6\linewidth}}{
        \cellcolor[rgb]{0.9,0.9,0.9}{
            \makecell[{{p{\linewidth}}}]{
                \texttt{\tiny{[GM$\langle$P2]}}
                \texttt{Sună bine să colaborăm sincer! } \\
\\ 
\texttt{Pentru mine: cartea valorează 1, geanta 2, iar șurubul 5. Deci șurubul îl apreciezi tu mai mult (8 vs 5), așa că e logic să{-}l iei tu. Cărțile nu au valoare pentru tine, dar au pentru mine, deci le iau eu. Geanta e la egalitate (2 pentru amândoi), așa că o iau eu ca să simplificăm.} \\
\\ 
\texttt{Propunere: tu iei șurubul, eu iau cele 3 cărți și geanta. Astfel maximizăm totalul (8 + 5 = 13). De acord?} \\
\\ 
\texttt{{[}Propunere: 3 cărți, 1 geantă{]}} \\
            }
        }
    }
     \\ \\

    \theutterance \stepcounter{utterance}  
    & & & \multicolumn{2}{p{0.3\linewidth}}{
        \cellcolor[rgb]{0.9,0.9,0.9}{
            \makecell[{{p{\linewidth}}}]{
                \texttt{\tiny{[GM$|$GM]}}
                \texttt{proposal} \\
            }
        }
    }
    & & \\ \\

    \theutterance \stepcounter{utterance}  
    & & & \multicolumn{2}{p{0.3\linewidth}}{
        \cellcolor[rgb]{0.9,0.9,0.9}{
            \makecell[{{p{\linewidth}}}]{
                \texttt{\tiny{[GM$|$GM]}}
                [3, 1, 0]
            }
        }
    }
    & & \\ \\

    \theutterance \stepcounter{utterance}  
    & & \multicolumn{4}{p{0.6\linewidth}}{
        \cellcolor[rgb]{0.9,0.9,0.9}{
            \makecell[{{p{\linewidth}}}]{
                \texttt{\tiny{[P1$\langle$GM]}}
                \texttt{Celălalt jucător a trimis o propunere secretă. Acum este rândul tău să trimiți o propunere proprie.} \\
            }
        }
    }
    & \\ \\

    \theutterance \stepcounter{utterance}  
    & \multicolumn{4}{p{0.6\linewidth}}{
        \cellcolor[rgb]{0.9,0.9,0.9}{
            \makecell[{{p{\linewidth}}}]{
                \texttt{\tiny{[P1$\rangle$GM]}}
                \texttt{{[}Propunere: 1 șurub, 1 geantă{]}} \\
            }
        }
    }
    & & \\ \\

    \theutterance \stepcounter{utterance}  
    & & & \multicolumn{2}{p{0.3\linewidth}}{
        \cellcolor[rgb]{0.9,0.9,0.9}{
            \makecell[{{p{\linewidth}}}]{
                \texttt{\tiny{[GM$|$GM]}}
                \texttt{proposal} \\
            }
        }
    }
    & & \\ \\

    \theutterance \stepcounter{utterance}  
    & & & \multicolumn{2}{p{0.3\linewidth}}{
        \cellcolor[rgb]{0.9,0.9,0.9}{
            \makecell[{{p{\linewidth}}}]{
                \texttt{\tiny{[GM$|$GM]}}
                [0, 1, 1]
            }
        }
    }
    & & \\ \\

    \theutterance \stepcounter{utterance}  
    & & & \multicolumn{2}{p{0.3\linewidth}}{
        \cellcolor[rgb]{0.9,0.9,0.9}{
            \makecell[{{p{\linewidth}}}]{
                \texttt{\tiny{[GM$|$GM]}}
                [[0, 1, 1], [3, 1, 0]]
            }
        }
    }
    & & \\ \\

\end{supertabular}
}

\end{document}
